% ============================================================================
% Borrador: metodología definitiva de MODELADO Y EVALUACIÓN (Etapa 2).
% Para pegar en el paper, después de la sección de datos/historial de Roy.
% Escrito en inglés como el resto del artículo. Todo lo que dice coincide con
% el código en src/hkr/ (revisar si se cambia el código).
% Las ecuaciones (2), (3) y (5) son las del paper de la Etapa 1.
% ============================================================================

\subsection{Model Ladder (Definitive)}

All models share one interface: they are fitted on the training years only
(1997--2003) and output one score per runner, where a higher score means a
higher chance of winning. Six rungs are trained, each adding a single idea to
the previous one, so that any improvement can be attributed to that idea:

\begin{enumerate}
  \item \textbf{Random within the field}: a constant score; its top-1 accuracy
        is reported as the exact expectation $\frac{1}{|R|}\sum_r 1/n_r$ rather
        than a single random draw.
  \item \textbf{Market favourite}: score $= \log(1/\text{odds})$ using the final
        win odds. It is an external reference only; odds are \emph{never} used
        as model features, so all learned models are odds-free and the market
        is the benchmark they are measured against.
  \item \textbf{Flat logistic regression} with $L_2$ regularisation, which also
        absorbs the collinearity between \texttt{horse\_rating} and
        \texttt{actual\_weight}. No class re-weighting is applied: the 8\% win
        rate is a structural property of races (one winner per field), not a
        sampling artefact.
  \item \textbf{Flat gradient boosting} (LightGBM, binary objective), the
        dominant approach in the reviewed literature.
  \item \textbf{Conditional logit} \cite{mcfadden}:
        $P(i \mid r) = \exp(\mathbf{x}_i^\top\boldsymbol\beta) /
        \sum_{j \in H_r} \exp(\mathbf{x}_j^\top\boldsymbol\beta)$, fitted by
        L-BFGS on the penalised race-level negative log-likelihood. It has no
        intercept (it cancels within a race) and isolates the effect of forcing
        probabilities to sum to one within each race. In dead heats the target
        mass is split equally among the joint winners.
  \item \textbf{LambdaMART} (LightGBM \texttt{lambdarank}), trained on the graded
        relevance label of Eq.~(5) with gain $2^{y}-1$, i.e.\ the same gain as
        Eq.~(2).
\end{enumerate}

\subsection{Preprocessing Inside the Pipeline}

Every transformation whose parameters are estimated from data is fitted on the
training split only and then applied unchanged to validation and test:

\begin{itemize}
  \item \emph{Linear models}: \texttt{prize} is imputed with the median of its
        \texttt{race\_class} (never a global median, since the missing values
        form a contiguous 1997--98 block), with an indicator
        \texttt{prize\_unrecorded}; remaining numeric features are
        median-imputed with missingness indicators and standardised;
        categorical features are one-hot encoded, grouping levels with fewer
        than 20 training rows and ignoring levels unseen in training.
  \item \emph{Tree models}: missing values are kept as \texttt{NaN} and handled
        natively by LightGBM (a debutant is not an average horse), and
        categorical levels are fixed to those observed in training.
\end{itemize}

Rows are sorted chronologically with each race contiguous, and the group-size
vector is validated against the row count before every ranker fit (Risk 2,
Table~II). A data contract check rejects any post-race column as a feature and
any race that spans two splits.

\subsection{Hyperparameter Search}

Hyperparameters are tuned with Optuna's TPE sampler (fixed seed) using
\emph{forward-chaining} cross-validation inside the training years: train on
1997--2000 and validate on 2001, then 1997--2001/2002, then 1997--2002/2003.
Folds are never shuffled, and neither 2004 nor 2005 is touched. The objective
is the mean MRR over the three folds, reported with its standard deviation. All
baselines with hyperparameters are tuned under the same protocol (15 trials for
the one-parameter linear models, 40 for each tree model), so the proposed model
receives no extra tuning advantage. For tree models the number of boosting
rounds is chosen by early stopping inside the folds and then frozen when the
final model is refitted on all training years. Search ranges, their rationale,
the boundary check, and the computational cost are documented in the
repository.

\subsection{Probability Calibration}

Scores are mapped to win probabilities with a race-wise softmax and a single
temperature $T$:
\begin{equation}
  \hat p_i = \frac{\exp(s_i/T)}{\sum_{j\in H_r}\exp(s_j/T)},
\end{equation}
where $T$ is fitted on the validation year (2004) by minimising the race-level
negative log-likelihood. Because $T$ is a single monotone parameter it leaves
every ranking metric unchanged; it only affects log loss, Brier score and
calibration error. Probabilistic metrics are reported exclusively on these
calibrated outputs.

\subsection{Evaluation Metrics}

Every metric is computed per race and then averaged over races, so that large
fields do not dominate:

\begin{itemize}
  \item \textbf{Top-1 accuracy}: the highest-scored runner won. In the 12 dead
        heats, picking any joint winner counts as a hit; tied scores count as
        the expected value of a random tie-break.
  \item \textbf{MRR}: mean reciprocal rank of the (best-ranked) winner.
  \item \textbf{NDCG@k}, $k\in\{1,3,6\}$, as in Eqs.~(2)--(3).
  \item \textbf{Within-race ROC-AUC} (winners vs.\ the rest of the field).
  \item \textbf{Race log loss} $-\log \hat p_{\text{winner}}$, the
        \textbf{race-level Brier score}
        $\sum_{i\in H_r}(\hat p_i - t_i)^2$, and the \textbf{expected
        calibration error} (ECE) over ten quantile bins, together with a
        reliability diagram.
\end{itemize}

Row-level accuracy is deliberately not reported: a classifier that never
predicts a winner already scores about 92\% (Section~III).

\textbf{Uncertainty.} With roughly 790 races in the validation year, the
standard error of top-1 accuracy is about 1.5 percentage points. Every
headline figure is therefore reported with a 95\% bootstrap confidence
interval obtained by resampling \emph{races} (2\,000 resamples). Differences
against the market favourite are computed as \emph{paired} differences on the
same resampled races; a difference is only claimed when its interval excludes
zero.

\textbf{Diagnostics.} Learning curves retrain each model on the last $k$
training years ($k=1,\dots,7$, always the most recent ones, so the distance to
2004 is constant) and compare training and validation MRR to diagnose
over- and under-fitting. Errors of the proposed model are broken down by field
size, by the winner's final odds, and by whether the model agreed with the
market favourite; the most confident misses and the worst-ranked winners are
inspected individually as boundary cases.

\subsection{Test Protocol}

The validation year (2004) is used for model comparison, temperature fitting
and the shrinkage constant $k$ of the historical features. The test year
(2005, 470 races) is evaluated exactly once, after all modelling decisions and
the explainability analysis are frozen; the evaluation script refuses to run a
second time. Because $k$ and $T$ were chosen on 2004, validation figures are
slightly optimistic for all models alike, whereas the test figures are not.

% ----------------------------------------------------------------------------
% NOTAS PARA EL EQUIPO (borrar antes de entregar)
% 1. Agregar a la bibliografía la referencia de McFadden (conditional logit):
%    D. McFadden, "Conditional logit analysis of qualitative choice behavior,"
%    in Frontiers in Econometrics, P. Zarembka, Ed. New York: Academic Press,
%    1974, pp. 105-142.  (\cite{mcfadden})
% 2. Cambios respecto a lo prometido en la Etapa 1 que hay que declarar:
%    - La Tabla II prometía reportar modelos "con y sin features de cuotas".
%      Se decidió NO usar cuotas como feature: el mercado es solo referencia.
%      Hay que actualizar esa fila de la Tabla II.
%    - La Ec. (4) prometía un promedio ponderado por recencia (decay). El
%      historial de v1 usa la tasa acumulada con shrinkage, sin decay. O se
%      implementa, o se dice explícitamente que se dejó fuera.
%    - Tabla V: "3 implausible weights" -> ahora se anula 1 (declared_weight
%      < 850 lb), coherente con el texto de la Sección V-D.
% ----------------------------------------------------------------------------
